% !TEX root = ../main.tex

\clearpage
\section{Benchmark and Evaluation Details}
\label{app:benchmarks}

Across all three benchmarks, tasks are evaluated in chronological order. The solver receives task $x_i$ with only the information available before its release time, while the evolver receives outcome labels only after the corresponding resolution time. Our analysis pipeline loads the first record for each \texttt{instance\_id}, so retries or duplicate logs do not change the reported metrics. The three benchmarks together exercise the three open-ended-stream deployment dimensions identified in \S\ref{sec:intro}.

\subsection{PolyBench}

\noindent\textbf{Source and composition.}
PolyBench is a Polymarket-derived prediction-market stream with 5{,}075 tasks from Feb 6--22, 2026. The stream spans politics, sports, finance, crypto, and entertainment markets, and each task is resolved against the official market outcome.

\noindent\textbf{Temporal ordering.}
For each market, we store the task release timestamp and the outcome resolution timestamp. Resolved labels are hidden from evolution until the corresponding market has resolved.

\noindent\textbf{Metrics.}
Let $N$ be the total number of markets in the stream and let $\mathcal{T}$ be the set of executed trades, excluding gated tasks, empty decisions, and \textsc{Skip} decisions. Let $s_i\in\{0,1\}$ indicate whether trade $i$ is correct, $b_i$ be its confidence-weighted investment, and $g_i$ be its realized profit. The metrics reported for PolyBench are:
\begin{align*}
\mathrm{Coverage} &= \tfrac{|\mathcal{T}|}{N}, \\
\mathrm{Acc}      &= 100 \cdot \tfrac{1}{N}\textstyle\sum_{i\in\mathcal{T}} s_i, \\
\mathrm{CWR}      &= 100 \cdot \tfrac{\sum_{i\in\mathcal{T}} g_i}{\sum_{i\in\mathcal{T}} b_i}, \\
\mathrm{Return}   &= \mathrm{Coverage} \cdot \mathrm{CWR}.
\end{align*}
Accuracy therefore rewards both broad coverage and correct decisions. Return is a portfolio-style profitability metric: CWR captures the dollar-weighted profit per unit invested over the markets the agent actually traded, and the Coverage scaling discounts a high CWR earned on only a thin slice of the stream. We report Return alongside Accuracy because Accuracy treats every trade equally and is blind to stake sizing, whereas under confidence-weighted investments a confident wrong trade can offset several confident correct ones; Return therefore reflects realized P\&L rather than mere directional correctness.

\subsection{CTF-Dojo}

\noindent\textbf{Source and composition.}
CTF-Dojo is a 261-challenge security-competition stream drawn from \texttt{pwncollege/ctf-archive}, chronologically ordered from 2011 to 2024. Challenges cover binary exploitation, web security, cryptography, reverse engineering, and forensics, exposing changes in challenge style and tooling over time.

\noindent\textbf{Sandbox and verification.}
Each challenge runs inside a per-task Docker sandbox with constrained network policy. Flags are submitted as text and verified by SHA-256 hash comparison against the official flag.

\noindent\textbf{Metrics.}
We report Pass@1, the percentage of challenges solved within the benchmark budget. All CTF-Dojo aggregate results use the same chronological ordering as the released stream.

\subsection{FutureX}

\noindent\textbf{Source and composition.}
FutureX is a 503-question event-forecasting stream over 82 days from Jan--Apr 2026, drawn from FutureX-Past. Questions cover finance, technology, geopolitics, and entertainment, with both English and Chinese-language variants. The zh-finance slice requires source discovery beyond default English-only retrieval.

\noindent\textbf{Temporal retrieval and filtering.}
FutureX tasks are historical, but web pages and search indices continue to change after the event. To avoid label leakage, each task is solved with a per-task cutoff date derived from its temporal metadata. In strict built-in retrieval, Wikipedia content is fetched through the revision API using the latest revision before the cutoff, DuckDuckGo results are fetched and filtered by extracted publication dates from \texttt{htmldate} and URL patterns, and structured economic series are queried with observation and realtime endpoints capped at the cutoff. For evolved tools executed through the sandbox, live command output is passed through an LLM temporal filter before the solver sees it: the filter receives the task, cutoff date, and retrieved content, then either returns \textsc{Clean} or replaces post-cutoff values, rows, and snippets with \texttt{[REDACTED]}. This preserves pre-cutoff evidence while blocking dynamic web content that would reveal the resolved answer.

\noindent\textbf{Metrics.}
We report Pass@1 under the official FutureX criterion. Slice-level human-steering analyses use the same pass criterion and aggregate tasks by the groups shown in Figure~\ref{fig:rq4_hitl_passrate}.

\section{Implementation and Reproducibility Details}
\label{app:experimental-details}

\noindent\textbf{Execution protocol.}
All main runs use provider-hosted LLM APIs with native tool calling and the same chronological task order used by the benchmarks. Reported numbers are generated from the corresponding \texttt{results.jsonl} files. Closed provider-hosted models do not expose parameter counts; we therefore report the evaluated task counts and evolution cycles rather than GPU-hours. The full-system runs contain 5{,}075/261/503 solve trajectories and 51/14/26 evolution cycles for PolyBench/CTF-Dojo/FutureX, respectively.

\noindent\textbf{Compute and resources.}
All model inference is performed through provider-hosted APIs; the primary Adaptive Auto-Harness runs use Claude Sonnet 4.6 for solving and Claude Opus 4.6 for evolution, with other provider-hosted models used only for the corresponding baseline rows. We do not train or fine-tune model weights, and no local GPU compute is used for model optimization. Local compute is used for orchestration, result aggregation, figure generation, and Docker-based benchmark execution, including the per-task CTF-Dojo and FutureX sandboxes. Because the evaluated closed models do not disclose parameter counts, task counts and evolution cycles are the main compute descriptors.

\noindent\textbf{Hyperparameters.}
Table~\ref{tab:hyperparams} lists the per-benchmark hyperparameters used by all main runs. Because temperature is zero throughout, the reported metrics are point estimates rather than averages over re-samples.

\begin{table}[t]
\centering\footnotesize
\setlength{\tabcolsep}{4pt}
\caption{Adaptive Auto-Harness hyperparameters across the three benchmarks. \emph{EGL} is the Expected-Gain-from-Learning trigger that gates whether a cycle runs. All runs use $T=0$ for both solver and evolver to attribute gains to the algorithm rather than sampling noise.}
\label{tab:hyperparams}
\resizebox{\columnwidth}{!}{%
\begin{tabular}{@{}lccc@{}}
\toprule
Hyperparameter & PolyBench & CTF-Dojo & FutureX \\
\midrule
\multicolumn{4}{@{}l}{\emph{Models}} \\
Solver & Sonnet 4.6 & Sonnet 4.6 & Sonnet 4.6 \\
Evolver & Opus 4.6 & Opus 4.6 & Opus 4.6 \\
Router & Sonnet 4.6 & Sonnet 4.6 & Sonnet 4.6 \\
\midrule
\multicolumn{4}{@{}l}{\emph{Sampling \& budget}} \\
Solver temperature & 0.0 & 0.0 & 0.0 \\
Evolver temperature & 0.0 & 0.0 & 0.0 \\
Solver max turns & 80 & 80 & 80 \\
Evolver max tokens & 128k & 128k & 128k \\
\midrule
\multicolumn{4}{@{}l}{\emph{Stream \& schedule}} \\
Total tasks & 5{,}075 & 261 & 503 \\
Batch size & 100 & 20 & 20 \\
Evolution cycles & 51 & 14 & 26 \\
EGL threshold & 0.05 & 0.05 & 0.05 \\
EGL window & 3 & 3 & 3 \\
Solve workers & 24 & 8 & 10 \\
\midrule
\multicolumn{4}{@{}l}{\emph{Multi-agent \& routing}} \\
Research parallel agents & 3 & 3 & 3 \\
Build/verify retries & 3 & 3 & 3 \\
Routing confidence threshold & 0.7 & 0.7 & 0.7 \\
\midrule
\multicolumn{4}{@{}l}{\emph{Sandbox}} \\
Solver sandbox network & none & none & bridge \\
Evolver sandbox network & none & none & bridge \\
\bottomrule
\end{tabular}}
\end{table}


\noindent\textbf{Seed harness.}
Table~\ref{tab:seed_inventory} reports what the evolver inherits before any cycle runs. The seed prompt is intentionally compact: PolyBench and CTF-Dojo seed prompts are 27 and 24 lines respectively, FutureX is longer because it documents the temporal-retrieval contract. All three benchmarks start with zero seed skills, tools, and memory entries, so any harness component beyond the seed prompt and the FutureX infrastructure scaffold is constructed by the evolver itself; this isolates the gains reported in the main results from prior hand-engineering of the seed.

\begin{table}[t]
\centering\footnotesize
\setlength{\tabcolsep}{6pt}
\caption{Seed harness shipped to the evolver before any cycle runs. \emph{Skills} counts top-level skill directories; \emph{Memory} counts JSONL entries; \emph{Infra} indicates whether the seed includes an infrastructure directory.}
\label{tab:seed_inventory}
\resizebox{\columnwidth}{!}{%
\begin{tabular}{@{}lrrrrl@{}}
\toprule
Benchmark & Prompt LOC & Skills & Tools & Memory & Infra \\
\midrule
PolyBench & 27 & 0 & 0 & 0 & no \\
CTF-Dojo & 24 & 0 & 0 & 0 & no \\
FutureX & 114 & 0 & 0 & 0 & yes \\
\bottomrule
\end{tabular}}
\end{table}


\noindent\textbf{Token cost and wall-clock.}
Table~\ref{tab:cost} summarises per-system token usage and wall-clock for the five most relevant systems. Solver tokens are summed from the per-task fields in \texttt{results.jsonl}; evolver-side tokens are omitted because the orchestrator did not persist them in the released artifacts. Wall-clock is the sum of per-task \texttt{elapsed} seconds and excludes orchestration overhead.

\begin{table}[t]
\centering\footnotesize
\setlength{\tabcolsep}{4pt}
\caption{Solver token cost and wall-clock per system. Tokens are summed from per-task \texttt{input\_tokens} / \texttt{output\_tokens} in \texttt{results.jsonl}. Wall-clock sums per-task \texttt{elapsed} seconds and excludes orchestration overhead. Evolver-side tokens were not persisted by the orchestrator in the released artifacts.}
\label{tab:cost}
\resizebox{\columnwidth}{!}{%
\begin{tabular}{@{}llrrrr@{}}
\toprule
System & Bench & In~(M~tok) & Out~(M~tok) & Hours & Tasks/h \\
\midrule
Sonnet (no-evo) & PolyBench & 35.3 & 5.0 & 25.6 & 198.3 \\
 & CTF-Dojo & 138.4 & 2.1 & 22.4 & 11.7 \\
 & FutureX & 55.5 & 0.7 & 34.2 & 14.7 \\
\midrule
A-Evolve & PolyBench & 365.4 & 4.2 & 23.9 & 212.3 \\
 & CTF-Dojo & 111.3 & 1.7 & 11.4 & 23.0 \\
 & FutureX & 268.8 & 1.2 & 13.8 & 36.5 \\
\midrule
Meta-Harness & PolyBench & 245.0 & 9.6 & 49.7 & 102.1 \\
 & CTF-Dojo & 142.7 & 2.2 & 20.0 & 13.0 \\
 & FutureX & 39.1 & 0.5 & 7.9 & 63.8 \\
\midrule
Multi-agent & PolyBench & 264.2 & 14.5 & 78.0 & 65.1 \\
 & CTF-Dojo & 180.2 & 2.5 & 21.2 & 12.3 \\
 & FutureX & 130.7 & 1.0 & 12.1 & 41.7 \\
\midrule
Full System & PolyBench & 233.2 & 12.2 & 59.5 & 85.2 \\
 & CTF-Dojo & 169.0 & 2.4 & 21.1 & 12.4 \\
 & FutureX & 25.6 & 0.5 & 6.6 & 75.7 \\
\bottomrule
\end{tabular}}
\end{table}



\noindent\textbf{Temporal reveal.}
Each task stores a release timestamp and, when available, a resolution timestamp. Solver calls are filtered against the release time. Evolution cycles receive the trajectory immediately but receive outcome feedback only after the task has resolved, so unresolved tasks remain unlabeled history rather than leaked supervision.

\noindent\textbf{Workspace artifacts.}
The evolver workspace persists a task board, research logs, verifier notes, tests, and architecture notes across cycles. These artifacts are separate from the solver workspace: the evolver may update harness files during evolution, while the solve-time router only inspects branch metadata and selects a branch for the incoming task.

\noindent\textbf{Branch replay protocol.}
For the routing analysis in Appendix~\ref{app:branch_perf}, every branch in the evolved harness tree is replayed on a curated task subset to construct the Oracle and Worst controls. Oracle and Worst are post-hoc diagnostic bounds; the deployed router sees only task context and branch metadata, not labels or branch outcomes.

\noindent\textbf{Human-steering records.}
Human-steering events are author-provided system interventions rather than recruited human-subject annotations. Each event is logged with the triggering phase, requested external signal, and workspace location where the response is recorded; Table~\ref{tab:hitl_log} reproduces the full event log. This makes steering auditable and keeps human input as source or access guidance rather than answer labels.

\section{Benchmark Non-Stationarity Diagnostics}
\label{app:nonstationarity}

Figures~\ref{fig:app_nonstationarity_polybench}--\ref{fig:app_nonstationarity_futurex} provide descriptive diagnostics for the temporal structure of the three streams. These plots are not used as evaluation metrics; instead, they show why the benchmarks are not static IID pools and why a harness fitted to earlier observations can become mismatched to later tasks. Most panels are computed from task metadata, task text, or benchmark-side properties; panels that use outcomes or a baseline solver are included only as descriptive solvability proxies.

\noindent\textbf{PolyBench.}
Prediction markets shift in both difficulty and tradability over the evaluated period (Figure~\ref{fig:app_nonstationarity_polybench}). Early markets are more often liquid and already decisive: the fraction of tradeable markets drops from 97\% early to 31\% late, and the fraction whose maximum price exceeds 0.95 drops from 44\% to 29\%. At the same time, near-even markets increase from 18\% to 35\%, meaning later tasks more often require evidence beyond simply following a strong market consensus. The market-price correctness proxy also changes over time, from 84\% early to 77\% late. Together, these shifts make a fixed prediction-market strategy brittle: calibration, abstention, and evidence gathering must adapt as the stream moves from liquid and decisive markets toward thinner and more ambiguous ones.

\noindent\textbf{CTF-Dojo.}
CTF-Dojo exposes a different form of non-stationarity: the stream expands into competitions and challenge conventions not present in the early history (Figure~\ref{fig:app_nonstationarity_ctf_dojo}). The cumulative number of source competitions keeps increasing across the chronological order, and by the late stream the fraction of tasks from competitions unseen in the first third reaches 100\%. The number of competitions represented in a 50-task window also varies substantially, so neighboring tasks can require different assumptions about file layout, scoring conventions, and intended exploitation style. The cross-competition score coefficient of variation further indicates that competitions are not interchangeable pools; each event can calibrate difficulty and challenge design differently. This motivates persistent construction of reusable security skills, but also cautions against treating early CTF experience as uniformly transferable.

\noindent\textbf{FutureX.}
FutureX shifts along source, language, and difficulty dimensions (Figure~\ref{fig:app_nonstationarity_futurex}). Batch-level baseline accuracy ranges from 20\% to 80\%, showing that chronological batches differ substantially in solvability. Later batches contain more Chinese-titled questions and more questions tied to platforms that are difficult to search directly, while the share of harder Level 3--4 questions increases sharply in the same region. Chinese-language answer requirements also appear mainly in later batches. These changes explain why FutureX stresses both construction and adaptation: the harness must acquire better source-finding and temporal retrieval behavior, while solve-time routing must select branches suited to the task's language, source, and difficulty profile.


\begin{figure}[t]
\centering
\includegraphics[width=\columnwidth]{figures/figure_c2_final.pdf}
\caption{\textbf{Evolver capability on CTF-Dojo.} Pass rate improves with stronger evolver models; higher budget helps Haiku and Sonnet but gives little additional gain once Opus already reaches high performance.}
\label{fig:app_c2_evolver_budget}
\end{figure}


\begin{figure}[t]
\centering
\includegraphics[width=\columnwidth]{figures/figure_c3_polybench_dilution.pdf}
\caption{\textbf{PolyBench workspace dilution.} A PolyBench-evolved workspace reaches the highest CWR, while combining all evolved workspaces sharply reduces CWR, supporting the need for specialized harness branches rather than a single dense harness.}
\label{fig:app_polybench_dilution}
\end{figure}



\begin{figure*}[t]
\centering
\includegraphics[width=0.9\linewidth]{figures/fig1_strategy_drift.pdf}
\caption{\textbf{PolyBench non-stationarity.} Market difficulty and tradability shift over time: later markets are less often decisive or liquid and more often near-even.}
\label{fig:app_nonstationarity_polybench}
\end{figure*}

\begin{figure*}[htb]
\centering
\includegraphics[width=0.8\linewidth]{figures/shift_crypto.pdf}
\caption{\textbf{CTF-Dojo non-stationarity.} The chronological stream keeps introducing new competitions and increases cross-competition variability, so early challenge experience is not uniformly transferable.}
\label{fig:app_nonstationarity_ctf_dojo}
\end{figure*}

\begin{figure*}[htb]
\centering
\includegraphics[width=0.8\linewidth]{figures/property_shift.pdf}
\caption{\textbf{FutureX non-stationarity.} Language, source accessibility, difficulty, and answer format shift across batches, creating solve-time harness mismatch when one static harness is reused for all tasks.}
\label{fig:app_nonstationarity_futurex}
\end{figure*}




\section{Further Experiments}
\label{app:further-experiments}

We include two supplemental diagnostics that clarify where the main gains come from without duplicating the RQ analyses in \S\ref{sec:experiments}.

\noindent\textbf{Evolver capability and construction budget.}
Figure~\ref{fig:app_c2_evolver_budget} varies the evolver model and construction budget on CTF-Dojo. Stronger evolvers achieve higher pass rates, while additional budget mainly helps weaker evolvers and saturates for the strongest model.



\noindent\textbf{Cross-domain workspace dilution.}
Figure~\ref{fig:app_polybench_dilution} compares PolyBench CWR when using workspaces evolved on different domains. The PolyBench-specific workspace performs best, while the all-evolved workspace loses 57 points of CWR, showing that mixing heterogeneous experience can dilute domain-relevant harness structure.


\section{Per-Domain and Per-Category Breakdowns}
\label{app:per_domain}

\noindent\textbf{CTF-Dojo.}
Table~\ref{tab:per_domain_ctf} reports Pass@1 by category. The Full System gains the most on \emph{web} ($+27$ over Sonnet) and \emph{crypto} ($+19$); \emph{binary/pwn} remains the hardest category at $14.8\%$ even after evolution and routing, consistent with the sandbox payload-handling bottleneck identified in \S\ref{sec:barriers}.

\begin{table}[t]
\centering\footnotesize
\setlength{\tabcolsep}{3pt}
\caption{Per-category Pass@1 (\%) on CTF-Dojo. Categories are parsed from the \texttt{detail} field; \emph{binary/pwn} merges the conventionally-equivalent CTF tags. Bold marks the best system per row.}
\label{tab:per_domain_ctf}
\resizebox{\columnwidth}{!}{%
\begin{tabular}{@{}lrrrrrrr@{}}
\toprule
Category & N & Sonnet & A-Evolve & Meta-H. & Multi & Adapt. & Full \\
\midrule
crypto & 74 & 52.7 & 66.1 & 56.2 & 67.6 & 55.6 & \textbf{72.1} \\
binary/pwn & 41 & 4.9 & 11.5 & 2.5 & 7.0 & 4.8 & \textbf{14.8} \\
web & 11 & 45.5 & 41.7 & 41.7 & 33.3 & 50.0 & \textbf{72.7} \\
reverse & 65 & 49.2 & 62.7 & 48.7 & 59.7 & 57.0 & \textbf{66.2} \\
forensics & 12 & 58.3 & 61.5 & 64.3 & 64.3 & \textbf{71.4} & 58.3 \\
misc & 58 & 20.7 & 23.5 & 29.5 & 34.1 & \textbf{40.5} & 25.6 \\
\midrule
Overall & 261 & 37.2 & 45.2 & 41.0 & 47.9 & 46.0 & \textbf{50.2} \\
\bottomrule
\end{tabular}}
\end{table}


\noindent\textbf{FutureX.}
Table~\ref{tab:per_domain_futurex} reports Pass@1 by language and inferred domain. English-language slices benefit most from evolution; the small \emph{zh-finance} slice illustrates the source-discovery bottleneck where neither evolution nor routing alone can recover when the platform is behind a search wall. Domains are inferred by keyword match on the question text, since FutureX results do not store the official domain tags; the \textit{other} bucket therefore aggregates unmatched questions.

\begin{table}[t]
\centering\footnotesize
\setlength{\tabcolsep}{3pt}
\caption{Per-slice Pass@1 (\%) on FutureX. Language is detected from Chinese characters in the question; domain is inferred by keyword match (\textit{other} catches unmatched questions). N is from the Sonnet baseline run. The \texttt{zh, geopolitics} row contains a single task and is reported for completeness only. Bold marks the best system per row.}
\label{tab:per_domain_futurex}
\resizebox{\columnwidth}{!}{%
\begin{tabular}{@{}llrrrrrrr@{}}
\toprule
Lang & Domain & N & Sonnet & A-Evolve & Meta-H. & Multi & Adapt. & Full \\
\midrule
en & finance & 76 & 21.1 & 53.9 & 21.1 & 52.6 & 43.4 & \textbf{56.6} \\
en & tech & 8 & 25.0 & \textbf{50.0} & 37.5 & \textbf{50.0} & 37.5 & \textbf{50.0} \\
en & geopolitics & 43 & 41.9 & 62.8 & 37.2 & 62.8 & 60.5 & \textbf{65.1} \\
en & sports & 62 & 38.7 & 58.1 & 37.1 & \textbf{59.7} & 56.5 & 54.8 \\
en & entertainment & 37 & 27.0 & 43.2 & 21.6 & 40.5 & \textbf{48.6} & 37.8 \\
en & other & 219 & 38.8 & 52.1 & 37.4 & \textbf{55.7} & 48.9 & 51.1 \\
\midrule
zh & finance & 10 & 0.0 & 0.0 & 0.0 & \textbf{30.0} & 0.0 & 20.0 \\
zh & geopolitics & 1 & 100 & 100 & 0 & 0 & 0 & 0 \\
zh & entertainment & 25 & 0.0 & 0.0 & 0.0 & 0.0 & 0.0 & 0.0 \\
zh & other & 22 & 0.0 & 0.0 & 0.0 & \textbf{4.5} & 0.0 & \textbf{4.5} \\
\midrule
& Overall & 503 & 31.0 & 47.5 & 29.4 & \textbf{49.5} & 44.1 & 47.3 \\
\bottomrule
\end{tabular}}
\end{table}


\noindent\textbf{PolyBench.}
Table~\ref{tab:per_domain_polybench} reports Accuracy/Return per inferred category. Sports dominates the portfolio ratio because liquid sports markets carry most of the dollar-weighted profit; politics carries near-zero return despite high accuracy. Categories are inferred by keyword match on the trajectory prompt's Event description.

\begin{table*}[t]
\centering\footnotesize
\setlength{\tabcolsep}{4pt}
\caption{Per-category PolyBench metrics. Categories are inferred by keyword match on the trajectory prompt's Event description; \textit{other} catches unmatched markets. Each cell reports \textbf{Accuracy~(\%) / Return~(\%)}. Bold marks the best system per row on Accuracy.}
\label{tab:per_domain_polybench}
\resizebox{\textwidth}{!}{%
\begin{tabular}{@{}lrcccccc@{}}
\toprule
Domain & N & Sonnet & A-Evolve & Meta-Harness & Multi-agent & Adaptive & Full System \\
\midrule
politics & 372 & 18.3/$-$5 & 11.3/$+$1 & 71.8/$-$2 & \textbf{87.9}/$-$3 & 87.6/$-$2 & 87.1/$-$4 \\
sports & 1{,}120 & 23.8/$+$3 & 16.4/$+$8 & 55.8/$+$594 & 81.7/$+$659 & 80.0/$+$642 & \textbf{84.6}/$+$596 \\
finance & 240 & 19.6/$+$2 & 18.3/$+$8 & 64.6/$+$7 & \textbf{81.7}/$+$10 & 77.1/$+$8 & 80.4/$+$9 \\
crypto & 447 & 20.4/$-$3 & 17.7/$+$3 & 66.9/$+$2 & 89.3/$+$4 & 88.4/$+$4 & \textbf{90.8}/$+$5 \\
entertainment & 218 & 22.0/$-$2 & 12.8/$+$1 & 64.7/$+$95 & \textbf{89.4}/$+$105 & 87.6/$+$101 & 89.0/$+$83 \\
other & 2{,}678 & 22.6/$+$3 & 20.7/$+$9 & 40.8/$+$356 & 75.4/$+$394 & 72.3/$+$409 & \textbf{76.2}/$+$374 \\
\midrule
Overall & 5{,}075 & 22.2/$+$2 & 18.4/$+$7 & 50.8/$+$320 & 79.8/$+$351 & 77.4/$+$352 & \textbf{80.9}/$+$330 \\
\bottomrule
\end{tabular}}
\end{table*}


\section{Multi-Agent Evolution Dynamics}
\label{app:evolution_dynamics}

Table~\ref{tab:evolution_dynamics} reports the full-stream system contrast (No-evo $\to$ Single-agent $\to$ Multi-agent) plus the timing of the multi-agent run, complementing rather than reproducing Figure~\ref{fig:rq2_stats}. The figure runs the four-phase evolver against \emph{No-memory} and \emph{No-feedback} ablations on a curated subset of samples; the table here instead reports the headline metric of each system on the full benchmark, matching the Multi-agent column of Table~\ref{tab:rq1_main}. The unique signal added is timing: the \emph{Peak} column gives the cycle index at which the multi-agent run's cumulative mean was highest, and a peak well before the final cycle is consistent with the overfitting trend in Figure~\ref{fig:c1}. On PolyBench the multi-agent peak occurs at cycle 22 of 51 and on FutureX at cycle 10 of 26, while CTF-Dojo continues to accumulate utility across all 14 cycles.

\begin{table}[t]
\centering\footnotesize
\setlength{\tabcolsep}{5pt}
\caption{Full-stream system contrast complementing the subset-based ablation in Figure~\ref{fig:rq2_stats}. \emph{No-evo} is the no-evolution Sonnet baseline; \emph{Single} is A-Evolve (single-agent evolver); \emph{Multi} is the four-phase multi-agent evolver. The \emph{Peak} column reports the cycle index at which the multi-agent run's cumulative mean was highest, and \emph{Cycles} is the total number of evolution cycles.}
\label{tab:evolution_dynamics}
\resizebox{\columnwidth}{!}{%
\begin{tabular}{@{}lrrrrrr@{}}
\toprule
Benchmark & Metric & No-evo & Single & Multi & Peak & Cycles \\
\midrule
PolyBench & Acc & 22.2 & 18.4 & \textbf{79.8} & 22 & 51 \\
CTF-Dojo & Pass@1 & 37.2 & 45.2 & \textbf{47.9} & 1 & 14 \\
FutureX & Pass@1 & 31.0 & 47.5 & \textbf{49.5} & 10 & 26 \\
\bottomrule
\end{tabular}}
\end{table}


\section{Routing Behaviour and Branch Performance}
\label{app:branch_perf}

This appendix expands the routing analysis behind Figure~\ref{fig:rq3_stats}. We use two related but distinct subsets per benchmark: (i)~the \emph{nav-run subset} (CTF-Dojo: 60 tasks, PolyBench: 100, FutureX: 80) on which the LLM router was deployed live and we observe its branch assignments; and (ii)~the \emph{replay subset} (CTF-Dojo: 40, PolyBench: 80, FutureX: 58), a subset of the same tasks for which every branch has been replayed end-to-end so each task carries a complete cross-venue score vector. The replay subset is necessarily smaller because some early-cycle branches did not exist for batch~1 tasks. Four per-task series are derived from the replay subset: \emph{Oracle} = best venue, \emph{Adapt} = the LLM router's actual choice on the nav run, \emph{Naive} = the fixed \texttt{main} venue (no branching), and \emph{Worst} = worst venue.

\noindent\textbf{Where the router sends each task.}
Table~\ref{tab:branch_perf} reports the per-branch routing volume and pass rate over the nav-run subset. The router prompt (Appendix~\ref{app:prompt-router}) does allow a \texttt{main} fallback when no branch matches strongly, but on these subsets the router always identifies a regime-specific branch and never invokes that fallback. Branches with low realised pass rates (e.g.\ \texttt{branch/pwn} on CTF-Dojo, \texttt{branch/lvl3} on FutureX) are not failing branches per se --- the router sends genuinely hard tasks to them, and the corresponding tasks have low Oracle pass rates on the replay subset as well (Table~\ref{tab:routing_stats}).

\begin{table}[t]
\centering\footnotesize
\setlength{\tabcolsep}{4pt}
\caption{Per-branch routing volume on the RQ4 navigation analysis subsets (\S\ref{sec:navloss}). Each row reports tasks the LLM router actually sent to that branch under the \emph{Adapt} condition, alongside the resulting Pass@1 (CTF-Dojo, FutureX) or HitRate among traded markets (PolyBench). On these subsets the router never invoked the \texttt{main} fallback that its prompt allows (Appendix~\ref{app:prompt-router}).}
\label{tab:branch_perf}
\resizebox{\columnwidth}{!}{%
\begin{tabular}{@{}llrr@{}}
\toprule
Bench & Branch & N routed & Pass / HitRate (\%) \\
\midrule
\multicolumn{4}{@{}l}{\emph{CTF-Dojo (60 tasks across 3 batches; per-task best-of-7 venues)}} \\
& \texttt{branch/crypto} & 20 & 50.0 \\
& \texttt{branch/rev} & 20 & 30.0 \\
& \texttt{branch/pwn} & 12 & 0.0 \\
& \texttt{branch/misc} & 4 & 75.0 \\
& \texttt{branch/web} & 2 & 0.0 \\
& \texttt{branch/forensics} & 2 & 50.0 \\
\midrule
\multicolumn{4}{@{}l}{\emph{PolyBench (100 tasks across 4 batches; per-task best-of-5 venues)}} \\
& \texttt{branch/sports} & 71 & 67.6 \\
& \texttt{branch/finance} & 14 & 21.4 \\
& \texttt{branch/culture} & 12 & 33.3 \\
& \texttt{branch/politics-world} & 3 & 66.7 \\
\midrule
\multicolumn{4}{@{}l}{\emph{FutureX (80 tasks across 3 batches; per-task best-of-5 venues)}} \\
& \texttt{branch/lvl1} & 28 & 53.6 \\
& \texttt{branch/lvl2} & 30 & 30.0 \\
& \texttt{branch/lvl3} & 8 & 0.0 \\
& \texttt{branch/lvl4} & 14 & 14.3 \\
\bottomrule
\end{tabular}}
\vskip -1em
\end{table}


\noindent\textbf{How much routing recovers of the adaptation gap.}
Table~\ref{tab:routing_stats} gives the headline Oracle/Adapt/Naive/Worst comparison on the replay subset, with 95\% bootstrap CIs and Holm--Bonferroni-corrected paired Wilcoxon \emph{p}-values. The Oracle$-$Naive gap, our empirical estimate of the adaptation loss $L_\mathrm{adapt}$, is large and significant on CTF-Dojo ($+37.5$\,pp, $p_\mathrm{adj}\!=\!4.8\times10^{-4}$) and PolyBench ($+8.8$\,pp CWR, $p_\mathrm{adj}\!=\!1.8\times10^{-5}$); on FutureX the gap is smaller and not significant after correction, consistent with the \S\ref{sec:barriers} finding that source acquisition rather than branch choice is the binding capability there. \emph{Adapt} closes a substantial fraction of the gap on CTF-Dojo and PolyBench but trails Naive slightly on FutureX, again reflecting the source-acquisition bottleneck.

\begin{table}[t]
\centering\footnotesize
\setlength{\tabcolsep}{4pt}
\caption{Numeric companion to Figure~\ref{fig:rq3_stats}. Replay-based Oracle/Adapt/Naive/Worst comparison on the RQ4 subsets (\S\ref{sec:navloss}). \emph{Oracle} is the best venue per task; \emph{Adapt} is the LLM router's choice; \emph{Naive} is the fixed \texttt{main} venue (no branching); \emph{Worst} is the worst venue per task. CTF-Dojo and FutureX use Pass@1 (\%); PolyBench uses CWR (\%). Means are reported with 95\% bootstrap CIs; gaps use the paired one-sided Wilcoxon signed-rank test with Holm--Bonferroni-corrected $p$-values. The \emph{Oracle$-$Naive} gap is the empirical adaptation loss $L_\mathrm{adapt}$.}
\label{tab:routing_stats}
\resizebox{\columnwidth}{!}{%
\begin{tabular}{@{}lccc@{}}
\toprule
 & CTF-Dojo & PolyBench & FutureX \\
 & (Pass\%) & (CWR\%) & (Pass\%) \\
\midrule
N tasks & 40 & 80 & 58 \\
\midrule
Oracle & 55.0\,[40.0, 70.0] & $+$12.0\,[$+$2.1, $+$21.2] & 46.6\,[32.8, 60.3] \\
Adapt  & 35.0\,[20.0, 50.0] & $+$5.9\,[$-$6.6, $+$17.4]  & 34.5\,[22.4, 46.6] \\
Naive  & 17.5\,[7.5, 30.0]  & $+$3.2\,[$-$7.5, $+$13.0]  & 39.7\,[27.6, 51.7] \\
Worst  & 7.5\,[0.0, 17.5]   & $-$10.6\,[$-$25.9, $+$4.0] & 22.4\,[12.1, 32.8] \\
\midrule
$L_\mathrm{adapt}$ (Oracle$-$Naive) & $+$37.5 & $+$8.8 & $+$6.9 \\
Adapt$-$Naive & $+$17.5 & $+$2.7 & $-$5.2 \\
\bottomrule
\end{tabular}}
\\[0.4em]
\end{table}


\noindent\textbf{Per-batch view.}
Table~\ref{tab:routing_per_batch} reports the same series broken out by batch on the replay subset (only batches with full cross-branch replays are shown, so e.g.\ CTF-Dojo includes batches 2 and 3 only). CTF-Dojo's branch tree improves between these two batches (Oracle $45\!\to\!65$\%) as new specialisations come online; PolyBench shows a quiet batch~4 followed by a sports-led recovery in batch~5; FutureX's batch~3 has the widest Oracle$-$Adapt headroom, where the router's level-based branches were less reliable than \texttt{main} on the same tasks.

\begin{table}[t]
\centering\footnotesize
\setlength{\tabcolsep}{4pt}
\caption{Per-batch breakdown of the Oracle/Adapt/Naive/Worst means on the RQ4 routing subsets. CTF-Dojo and FutureX use Pass@1 (\%); PolyBench uses CWR (\%). Adapt closes the Oracle$-$Naive gap most consistently on CTF-Dojo, where branch quality is stable; on FutureX the gap is small and noisy because source acquisition rather than branch choice dominates the failure mode.}
\label{tab:routing_per_batch}
\resizebox{\columnwidth}{!}{%
\begin{tabular}{@{}lrrrrrr@{}}
\toprule
Benchmark & Batch & N & Oracle & Adapt & Naive & Worst \\
\midrule
\multirow{2}{*}{CTF-Dojo (Pass\%)}
 & 2 & 20 & 45.0 & 30.0 & 15.0 & 10.0 \\
 & 3 & 20 & 65.0 & 40.0 & 20.0 & 5.0 \\
\midrule
\multirow{4}{*}{PolyBench (CWR\%)}
 & 2 & 20 & $+$16.7 & $+$11.6 & $+$7.8 & $-$100.0 \\
 & 3 & 20 & $+$13.0 & $+$7.8 & $+$1.9 & $-$8.9 \\
 & 4 & 20 & $+$1.1 & $+$1.1 & $+$1.4 & $-$5.1 \\
 & 5 & 20 & $+$14.7 & $+$5.2 & $+$0.5 & $-$4.7 \\
\midrule
\multirow{3}{*}{FutureX (Pass\%)}
 & 2 & 20 & 45.0 & 40.0 & 45.0 & 40.0 \\
 & 3 & 19 & 57.9 & 26.3 & 52.6 & 21.1 \\
 & 4 & 19 & 36.8 & 36.8 & 21.1 & 5.3 \\
\bottomrule
\end{tabular}}
\vskip -1em
\end{table}


\section{Human-in-the-Loop Event Log}
\label{app:hitl_log}

Table~\ref{tab:hitl_log} reports the full HITL event log from the FutureX RQ5 run (\S\ref{sec:hitl}). The run uses the curated 5-batch FutureX stream (Design~C: 100 tasks, 5\,$\times$\,20) with two engineered regime shifts, a Sonnet 4.6 solver, an Opus 4.6 evolver, and \texttt{hitl\_enabled=true}. The Analyst and Builder decide when to invoke each hook; the human responds via Telegram from a pre-authored cheat-sheet. Two P2 (research-phase) events fire at cycle 1 to bootstrap the search pipeline, and one substantive P3 (task-board) event fires at cycle 3 to steer the evolver toward Western and Chinese specialty endpoints. The remaining P3 prompts (cycles 1, 2, 4, and 5) return \texttt{skip}, matching the cheat-sheet protocol that the human only intervenes when the cheat-sheet has a relevant entry. The slice-level lift on each regime ($0$, $+5$, $+20$, $+15$, $0$ on regimes 1–5) is reported in Figure~\ref{fig:rq4_hitl_passrate}.

\begin{table*}[t]
\centering\footnotesize
\setlength{\tabcolsep}{4pt}
\caption{Complete human-in-the-loop event log from the FutureX RQ5 run (Design~C, 5 batches $\times$ 20 tasks). \emph{P2 (cred.)} is the research-phase credential hook; \emph{P3 (board)} is the task-board steering hook. Two P2 events fire at cycle~1 to bootstrap the search pipeline; one substantive P3 event fires at cycle~3 to direct the evolver toward Western and Chinese specialty endpoints. The remaining P3 prompts return \texttt{skip}, matching the cheat-sheet protocol. API tokens supplied via Telegram are redacted.}
\label{tab:hitl_log}
\renewcommand{\arraystretch}{1.15}
\begin{tabular}{@{}rllp{0.22\textwidth}p{0.34\textwidth}@{}}
\toprule
Cycle & Hook & Key & Trigger context & Human response \\
\midrule
1 & P2 (cred.) & \texttt{EXA\_API\_KEY} & Phase-2 research needs Exa search API. & \texttt{[REDACTED]} (key supplied) \\
1 & P2 (cred.) & \texttt{SERPER\_API\_KEY} & Phase-2 research needs Serper Google API. & \texttt{[REDACTED]} (key supplied) \\
1 & P3 (board) & cycle1 & 7 tasks fail: solver has no web search tool. & skip \\
2 & P3 (board) & cycle2 & 9 tasks fail: search-pipeline entry never invokes search functions. & skip \\
3 & P3 (board) & cycle3 & deterministic fallback used; structured-data API gap. & ``Specialty data tasks need direct endpoint integrations beyond generic web search. Build skills/tools for Western (US equities OHLC, Box Office Mojo, $\ldots$) and Chinese (Eastmoney secid, Maoyan film board, KolRank, $\ldots$) endpoints.'' \\
4 & P3 (board) & cycle4 & 17 tasks fail: deterministic fallback used. & skip \\
5 & P3 (board) & cycle5 & 20 tasks fail: Chinese niche-ranking API absent. & skip \\
\bottomrule
\end{tabular}
\end{table*}


\section{Run-Detail Analysis}
\label{app:run_details}

Table~\ref{tab:run_turns_time} reports per-task turn counts and elapsed seconds, complementing the aggregate cost summary in Table~\ref{tab:cost}. Both distributions are right-skewed on CTF-Dojo and FutureX, where a small number of long-running tasks pull the mean above the median; we therefore report both. CTF-Dojo's solver budget is consistently saturated --- Sonnet's mean of 89.4 turns reflects a large mass of tasks that loop until the cap, whereas the evolved variants typically reach a flag (or give up) earlier. PolyBench is dominated by direct decisions: most tasks submit immediately (turn count 1) because the prompt already contains the market context the solver needs.

\begin{table*}[t]
\centering\footnotesize
\setlength{\tabcolsep}{4pt}
\caption{Per-task solver turns and wall-clock seconds across systems and benchmarks. A \emph{turn} is one tool call, including the final \texttt{submit}; a task that submits directly without other tool use therefore counts as 1. \emph{$\overline{\mathrm{turns}}$} and \emph{$\overline{\mathrm{sec}}$} are arithmetic means; \emph{median} columns are added because both distributions are right-skewed on CTF-Dojo and FutureX. Wall-clock excludes orchestration overhead.}
\label{tab:run_turns_time}
\resizebox{\textwidth}{!}{%
\begin{tabular}{@{}l rrrr rrrr rrrr@{}}
\toprule
& \multicolumn{4}{c}{PolyBench} & \multicolumn{4}{c}{CTF-Dojo} & \multicolumn{4}{c}{FutureX} \\
\cmidrule(lr){2-5} \cmidrule(lr){6-9} \cmidrule(lr){10-13}
System & $\overline{\mathrm{turns}}$ & med & $\overline{\mathrm{sec}}$ & med & $\overline{\mathrm{turns}}$ & med & $\overline{\mathrm{sec}}$ & med & $\overline{\mathrm{turns}}$ & med & $\overline{\mathrm{sec}}$ & med \\
\midrule
Sonnet & 1.0 & 1.0 & 18.2 & 18.1 & 89.4 & 48.0 & 308.4 & 282.8 & 13.4 & 16.0 & 244.4 & 271.4 \\
A-Evolve & 1.7 & 2.0 & 17.0 & 15.7 & 17.8 & 10.0 & 156.7 & 76.9 & 16.9 & 6.0 & 98.8 & 41.9 \\
GEPA & 2.1 & 2.0 & 30.4 & 22.5 & 50.1 & 33.0 & 250.3 & 219.7 & 0.8 & 1.0 & 4.6 & 4.6 \\
Meta-Harness & 2.4 & 2.0 & 35.2 & 29.5 & 51.3 & 39.0 & 275.9 & 244.4 & 4.6 & 4.0 & 56.4 & 50.4 \\
Continual H. & 2.3 & 2.0 & 27.4 & 21.6 & 15.3 & 12.0 & 115.3 & 82.3 & 6.8 & 5.0 & 167.8 & 112.7 \\
SkillOS & 2.7 & 2.0 & 32.1 & 23.0 & 31.4 & 16.0 & 224.2 & 161.0 & 5.3 & 4.0 & 120.3 & 84.1 \\
OctoTools & 5.2 & 4.0 & 68.6 & 63.5 & 44.1 & 35.0 & 264.8 & 253.7 & 1.0 & 1.0 & 10.7 & 10.5 \\
Multi-agent & 5.7 & 6.0 & 55.3 & 48.7 & 40.1 & 35.0 & 293.1 & 245.2 & 11.3 & 3.0 & 86.4 & 36.2 \\
Adaptive & 2.8 & 3.0 & 43.7 & 38.3 & 29.7 & 25.0 & 281.9 & 238.1 & 10.1 & 8.0 & 74.5 & 55.7 \\
Full System & 5.4 & 6.0 & 42.2 & 40.5 & 36.4 & 42.0 & 290.9 & 255.9 & 4.2 & 4.0 & 47.6 & 38.6 \\
\bottomrule
\end{tabular}}
\end{table*}


\clearpage
\onecolumn
\section{System Prompts}
\label{app:prompts}

We reproduce the verbatim system prompts used by Adaptive Auto-Harness, exactly as the agents receive them. Curly-brace placeholders such as \texttt{\{benchmark\_context\}}, \texttt{\{regime\}}, \texttt{\{workspace\_extras\}}, and \texttt{\{categories\}} are substituted by the framework at runtime per benchmark or per regime; we leave them in place so the templating is visible.

\label{app:prompt-router}
\begin{promptlisting}[Solve-time router]
You are a fast task router. For each task you must pick exactly one branch
from the list provided.

Decision procedure:
1. Read each branch's README, focusing on its "When to route here" section.
2. Match the task description against each branch's positive signals
   (categories, keywords, metadata) and negative signals (NOT applicable).
3. If exactly one specialized branch matches strongly → route there with
   high confidence (>= 0.7).
4. If multiple branches partially match or none match clearly → route to
   "main" with confidence reflecting your uncertainty.
5. Branches without a clear "When to route here" section should be treated
   as low-confidence candidates — prefer main unless the regime is obvious.

Output JSON only:
{"branch": "main", "confidence": 0.8, "reason": "brief reason"}
\end{promptlisting}

\label{app:prompt-researcher}
\begin{promptlisting}[Researcher]
You are a research agent — PHASE 2 of 4 in the evolution cycle.

PHASE SEQUENCE:
  1. ANALYZE  → analyst identified gaps, wrote task_board.md
  2. RESEARCH (you) → discover solutions for your assigned regime
  3. BUILD    → builder reads YOUR research to write code
  4. VERIFY   → verifier tests what the builder created

UPSTREAM: The analyst assigned you regime "{regime}" because tasks
are failing in that area.

DOWNSTREAM: The builder will read your research records to decide
what code to write. For each approach you test, document:
- How to implement it (endpoint, arguments, response format)
- What it covers and what it doesn't
- Whether it works reliably (tested with real calls)

RESEARCH APPROACH:
1. Read /solver_workspace/ to see what's already implemented
2. Read /evolver_workspace/research_log.jsonl to avoid retesting
3. Search the web for solutions, libraries, APIs, reference code
4. Test NEW approaches that complement what already exists

You can utilize Full network access if that's available.

Write findings to /evolver_workspace/tests/research_{regime}.jsonl.

WORKSPACE LAYOUT:
  /solver_workspace/    — solver workspace (read current code)
  /evolver_workspace/   — evolution state
    research_log.jsonl  — existing records (read to avoid retesting)
    tests/              — write findings HERE
  /trajectories/        — READ-ONLY solver conversations per task

BASH OUTPUT IS CAPPED AT 100 KB PER CALL (first 50 KB + last 50 KB,
middle elided). When reading trajectories or large logs, prefer
`jq`, `grep`, `head`, `tail` over raw `cat` to keep context focused.

{benchmark_context}
\end{promptlisting}

\label{app:prompt-builder}
\begin{promptlisting}[Builder]
You are a builder — PHASE 3 of 4 in the evolution cycle.

PHASE SEQUENCE:
  1. ANALYZE  → analyst identified failure patterns in task_board.md
  2. RESEARCH → agents discovered solutions in research_log.jsonl
  3. BUILD (you) → implement solutions from verified research
  4. VERIFY   → verifier tests YOUR code

UPSTREAM: Read research_log.jsonl for verified approaches (works=true).
The task_board tells you which capabilities matter most.
The architecture.md shows what's already built.

DOWNSTREAM: The verifier will test your code. If verification fails,
you get the report and can retry (max 3 attempts).

WORKSPACE LAYOUT:
  /solver_workspace/          — the solver's workspace
    prompts/system.md         — solver prompt (update if needed)
    tools/                    — evolved tool scripts
    infra/                    — evolved infrastructure (if applicable)
{workspace_extras}  /evolver_workspace/
    task_board.md             — failure regimes from analyst
    research_log.jsonl        — verified approaches from research
    architecture.md           — UPDATE with what you built
  /trajectories/              — READ-ONLY solver conversations

Read existing code first. Extend, don't rewrite.
Do NOT run git — the framework handles commits.

BASH OUTPUT IS CAPPED AT 100 KB PER CALL (first 50 KB + last 50 KB,
middle elided). When scanning large files, prefer `jq`, `grep`,
`head`, `tail` over raw `cat`.

{benchmark_context}
\end{promptlisting}

\label{app:prompt-verifier}
\begin{promptlisting}[Verifier]
You are a verification agent — PHASE 4 of 4 in the evolution cycle.

PHASE SEQUENCE:
  1. ANALYZE  → analyst identified failure patterns
  2. RESEARCH → agents discovered solutions
  3. BUILD    → builder wrote code implementing those solutions
  4. VERIFY (you) → test the builder's code before it goes live

UPSTREAM: The builder just wrote/modified code in /solver_workspace/.
Your job is to verify it works before the next batch.

DOWNSTREAM: If PASS, the code goes live. If FAIL, the builder gets
your report and retries. Be specific about what failed.

WORKSPACE LAYOUT:
  /solver_workspace/    — solver workspace (test code here)
  /evolver_workspace/   — evolution state
    task_board.md       — context for test cases
    tests/              — write your test scripts HERE
  /trajectories/        — READ-ONLY solver conversations

Output: VERDICT: PASS or FAIL, then list each test with its result.

BASH OUTPUT IS CAPPED AT 100 KB PER CALL (first 50 KB + last 50 KB,
middle elided). Prefer `jq`, `grep`, `head`, `tail` over raw `cat`
when inspecting large files.

{benchmark_context}
\end{promptlisting}

\label{app:prompt-analyst}
\begin{promptlisting}[Analyst]
You are a failure analyst — PHASE 1 of 4 in the evolution cycle.

PHASE SEQUENCE:
  1. ANALYZE (you) → write task_board.md with failure patterns + priorities + TARGET
  2. RESEARCH      → agents investigate top-K gaps from YOUR task board
  3. BUILD         → builder implements solutions PER TARGET (main first, then branches)
  4. VERIFY        → verifier tests what the builder created per target

Your task board DIRECTLY drives what gets researched and built next.
Be specific about what capability is missing — vague gaps lead to
unfocused research.

This system uses NAVIGATION: git branches isolate solver strategies.
Each task is routed to the best branch before solving. Your job includes
deciding WHERE each fix should land.

CRITICAL: TRANSFERABILITY ANALYSIS (do this BEFORE writing the task board)

Naive evolution accumulates "shortcut artifacts" — skills, tools, prompt rules,
or memory entries that helped one batch but break later tasks with different
distribution. Examples seen in past runs:
  - A "stop searching at 6 queries" rule helped batch 7 (sports) at 95% but
    crashed batches 14-20 (Chinese music) to 15% because those tasks needed
    longer searches.
  - A `bls_forgery.md` skill with 40s overhead and hard-coded BLS curve
    parameters helped 1 BLS task but slowed all other crypto tasks.
  - 22 hard-coded "wrong flag" entries in system.md memorized past failures
    without abstracting the underlying lesson, bloating the prompt.

Your job in Phase 1 is to AUDIT the evolution state for non-transferable
artifacts BEFORE proposing more fixes. Use bash to:

  1. Compare per-CATEGORY pass rates ACROSS cycles (read multiple
     trajectories/batch_NNNN/index.txt files). If category X had 80% in
     early batches and 30% in recent batches — that's degradation.
  2. List recent additions to skills/, tools/, prompts/system.md, memory/
     and ask: which CATEGORIES does each artifact help vs. hurt?
     - Look for hard-coded task IDs, year-specific logic, single-domain rules
     - Look for prompt rules added recently that contradict older rules
  3. Check the strategy_tree.md routing stats: if a branch is dragging
     overall performance down (worse than main on its routed tasks),
     mark it for retirement.

For each non-transferable artifact found, add a "## Toxic Artifacts" section
to the task board listing:
  - artifact_name: helps {categories} but hurts {categories} → ACTION:
    move-to-branch/<name>, deprecate, or rewrite-as-general

BRANCHING DECISION (TARGET per regime):
- TARGET: main — the fix is domain-generalizable AND verified to NOT degrade
  any previously-passing category
- TARGET: branch/<name> — the fix is regime-specific OR has been observed to
  help one category while hurting another (isolate it from main)
- TARGET: branch/<existing-name> — an existing branch already handles this regime
- Do NOT create new branches for < 2 tasks or single-cycle observations
- Prefer branch isolation when in doubt: keeping main clean is more important
  than maximizing main's per-batch peak

WORKSPACE LAYOUT:
  /solver_workspace/          — the solver's workspace (may be on any branch)
  /evolver_workspace/         — shared evolution state
    task_board.md             — YOUR OUTPUT
    research_log.jsonl        — what's been researched so far
    architecture.md           — what's been built so far
    strategy_tree.md          — current branch descriptions + routing stats
    evolution/observations/   — batch results (revealed feedback only)
  /trajectories/              — READ-ONLY per-task solver conversations

USE BASH to deeply analyze:
- /trajectories/batch_NNNN/index.txt for per-task category/year/outcome summary
  (read this FIRST to see regime distribution before reading full trajectories)
- /trajectories/ for full solver conversations per task
- /evolver_workspace/strategy_tree.md for branch performance
- /evolver_workspace/evolution/observations/ for batch results
- /solver_workspace/ to see current evolved code

BASH OUTPUT IS CAPPED AT 100 KB PER CALL (first 50 KB + last 50 KB,
middle elided). Trajectory JSONs on security/crypto tasks can be
200+ KB each — `cat` will truncate them, and reading several in a row
will exhaust your context. PREFER:
  - `jq '.steps[].tool_use // .steps[].output' traj.json | head -200` to
    see tool calls/outputs without the raw conversation bulk
  - `jq -r '.steps[-5:]' traj.json` to inspect only the final few steps
  - `grep -n ERROR|FAIL|flag traj.json` to locate specific signals
  - `ls -lS /trajectories/ | head` to find the largest trajectories first
  - `wc -l traj.json` before `cat` to check size
Reserve raw `cat` for files under ~50 KB.

NAVIGATION CONTEXT (provided in the user prompt):
- Strategy tree: existing branches and their per-task routing performance
- Routing summary: which tasks went to which branch this batch, and passed/failed

PRIVACY: feedback_archive.jsonl is masked. The observations/ files
contain all feedback you are allowed to see under temporal-reveal.

OUTPUT FORMAT: task board in the exact markdown format with TARGET annotations:
```
## Failure Patterns (Cycle N)
- regime_name: N tasks ... PRIORITY: HIGH → TARGET: main
- other_regime: M tasks ... PRIORITY: MEDIUM → TARGET: branch/regime-name
```

No conversational text in the final output.

{benchmark_context}
\end{promptlisting}

